\chapter{Wireless Principles}\label{ch:wireless}
\bp{1.5}

\begin{outcomes}
By the end of this chapter you will be able to:
\begin{tight}
  \item \textbf{Select} a band and a channel plan, and explain why 2.4\,GHz has
        only three usable channels.
  \item \textbf{Describe} the RF characteristics that decide whether a link
        works: signal strength, noise, SNR, and what happens to radio waves on
        the way.
  \item \textbf{Compare} WEP, WPA, WPA2 and WPA3, and distinguish personal from
        enterprise authentication.
  \item \textbf{Identify} causes of interference, and tell co-channel
        interference apart from adjacent-channel interference.
\end{tight}
\end{outcomes}

\begin{prereq}
\chref{ethernet} --- wireless is a shared medium with collision \emph{avoidance}
rather than detection, and the contrast is the quickest way in.
\chref{edgehost} for how access points attach to the wired network.
\chref{aaa} for RADIUS, which is what ``enterprise'' wireless authenticates
against.
\end{prereq}

\begin{keyterms}
band $\bullet$ channel $\bullet$ non-overlapping $\bullet$ channel bonding
$\bullet$ dBm $\bullet$ RSSI $\bullet$ SNR $\bullet$ noise floor $\bullet$
attenuation $\bullet$ SSID $\bullet$ BSSID $\bullet$ WPA2 $\bullet$ WPA3
$\bullet$ SAE $\bullet$ 802.1X $\bullet$ co-channel interference
\end{keyterms}

\begin{examnote}[Four named sub-topics, and the verb is \emph{describe}]
Topic~\tp{1.5} is \emph{describe wireless principles}, with four sub-topics named
explicitly: \tp{1.5.a} band and channel selection, \tp{1.5.b} RF characteristics,
\tp{1.5.c} security protocols, \tp{1.5.d} cause of interference.

There is \textbf{no wireless configuration on this blueprint.} v1.1 asked you to
build a WLAN in a controller GUI; v2.0 does not. Do not spend time learning WLC
menus for this exam --- spend it on the four lists in this chapter, because that
is exactly what is examined.

Access points do appear elsewhere: \tp{2.2} covers connecting them to a switch
port, which is Chapter~9's material, and \tp{1.6} covers the client end, which is
Chapter~28's.
\end{examnote}

\section{Why wireless is different}

A switch port is a private, full-duplex conversation between two devices. A radio
channel is a room everybody shouts across. Three consequences shape everything
else in this chapter.

\begin{tight}
  \item \textbf{It is half duplex and shared.} Only one device on a channel may
        transmit at a time. Ten clients on one access point share the airtime
        between them --- the advertised data rate is not per-client bandwidth.
  \item \textbf{Collisions cannot be detected, only avoided.} A radio cannot
        listen while it transmits, so Ethernet's CSMA/CD is impossible. Wireless
        uses CSMA/CA: listen, wait a random interval, transmit, and expect an
        acknowledgement for every frame.
  \item \textbf{The medium is unlicensed and shared with strangers.} Your
        neighbour's network, a microwave oven and a Bluetooth headset all have
        the same right to the spectrum that you do.
\end{tight}

\begin{definitionbox}[Four names that get confused]
\begin{tight}
  \item \textbf{SSID} --- the network's name, as a human sees it. Several APs
        normally share one.
  \item \textbf{BSSID} --- the MAC address of one radio on one AP. This is what a
        client is actually associated \emph{to}.
  \item \textbf{BSS} --- one AP's coverage area. \textbf{ESS} --- several APs
        sharing an SSID so a client can roam between them.
  \item \textbf{Roaming} --- a client deciding, on its own, to move from one
        BSSID to another within the same ESS. \emph{The client decides}, not the
        network. This matters in Chapter~28.
\end{tight}
\end{definitionbox}

\section{Bands and channels}

\begin{longtable}{@{}L{2.4cm}L{3.4cm}L{4.2cm}L{4.9cm}@{}}
\toprule
\rowcolor{primary!15}
\textbf{Band} & \textbf{Usable width} & \textbf{Non-overlapping 20\,MHz channels} & \textbf{Character} \\
\midrule
\endhead
2.4\,GHz & About 83\,MHz &
\textbf{Three}: 1, 6 and 11 &
Long range, penetrates walls, hopelessly crowded \\
\rowcolor{lightbg}
5\,GHz & About 500\,MHz &
Around \textbf{25}, some requiring DFS &
Shorter range, far more room, the sensible default \\
6\,GHz & 1200\,MHz &
\textbf{59} &
Wi-Fi 6E and later only. Clean, but short range and newer clients only \\
\bottomrule
\end{longtable}

\begin{alertbox}[Why 2.4\,GHz has only three channels]
The band is divided into channels numbered 1 to 11 in North America (1 to 13 in
much of the rest of the world), spaced 5\,MHz apart. But each channel is
\textbf{20\,MHz wide}. A 20\,MHz signal spaced 5\,MHz from its neighbour overlaps
it heavily.

To avoid overlap you need five channel numbers of separation --- which in a band
of eleven gives you exactly three: \textbf{1, 6 and 11}. Any other plan puts two
networks partly on top of each other, which is worse than putting them fully on
top of each other. That last point is the one the exam tests, and it is explained
below.
\end{alertbox}

\ccnafig{%
\begin{tikzpicture}[font=\scriptsize]
  \draw[neutral!60, line width=0.8pt] (0,0) -- (13.2,0);
  \foreach \c/\x in {1/0.6, 2/1.2, 3/1.8, 4/2.4, 5/3.0, 6/3.6, 7/4.2, 8/4.8,
                     9/5.4, 10/6.0, 11/6.6} {
    \draw[neutral!60, line width=0.6pt] (\x,0) -- (\x,-0.14);
    \node[font=\tiny, text=neutral!85!black] at (\x,-0.38) {\c};
  }
  \node[font=\tiny, text=neutral!85!black, anchor=west] at (7.1,-0.38)
      {channel number, 5\,MHz apart};

  \foreach \x/\col in {0.6/dom2, 3.6/dom3, 6.6/dom4} {
    \draw[\col, fill=\col!14, line width=0.9pt]
      (\x-1.2,0.15) -- (\x-0.6,1.25) -- (\x+0.6,1.25) -- (\x+1.2,0.15) -- cycle;
  }
  \node[font=\tiny\bfseries, text=dom2!45!black] at (0.6,0.72) {1};
  \node[font=\tiny\bfseries, text=dom3!45!black] at (3.6,0.72) {6};
  \node[font=\tiny\bfseries, text=dom4!45!black] at (6.6,0.72) {11};

  % an overlapping choice
  \draw[accent, dash pattern=on 2.5pt off 2pt, fill=accent!10, line width=0.9pt]
    (1.2,0.15) -- (1.8,1.25) -- (3.0,1.25) -- (3.6,0.15) -- cycle;
  \node[font=\tiny\bfseries, text=accent!70!black] at (2.4,0.72) {3};
  \node[font=\tiny, text=accent!70!black, align=center, anchor=north]
    at (2.4,-0.75) {channel 3 overlaps\\both 1 and 6};

  \node[font=\tiny, text=neutral!85!black, anchor=west] at (8.4,0.9)
      {each channel is 20\,MHz wide};
\end{tikzpicture}}%
{The 2.4\,GHz band. Channels 1, 6 and 11 are the only trio that does not overlap;
anything between them interferes with two networks instead of sharing one.}{chan24}

\subsection{Channel bonding}

Wider channels carry more data. 802.11n and later can bond adjacent 20\,MHz
channels into 40, 80 or 160\,MHz.

\begin{alertbox}[The trade-off, which is examinable]
Every doubling of channel width roughly doubles throughput --- and \textbf{halves
the number of non-overlapping channels available}, and raises the noise floor by
about 3\,dB because the receiver is listening to twice as much spectrum.

So: bonding is sensible at 5 and 6\,GHz where there is room, and in a
low-density deployment. It is \textbf{never} sensible at 2.4\,GHz, where a
40\,MHz channel consumes most of the band and leaves no clean plan for anyone,
including you.
\end{alertbox}

\begin{conceptbox}[DFS channels]
Part of the 5\,GHz band is shared with weather and military radar. Access points
using those channels must run \term{dynamic frequency selection}: listen before
transmitting, and vacate the channel within seconds if radar appears.

The practical consequence is a real fault you will meet: a DFS event moves an AP
to a different channel and every client on it is briefly disconnected. If a site
reports unexplained drops on 5\,GHz at irregular intervals --- particularly near
an airport or a coast --- DFS is the first thing to check.
\end{conceptbox}

\section{RF characteristics}

\subsection{Measuring power in decibels}

Radio power spans an enormous range, so it is measured logarithmically.

\begin{tight}
  \item \textbf{dBm} --- power relative to 1\,milliwatt. 0\,dBm is 1\,mW.
        Absolute.
  \item \textbf{dB} --- a ratio, with no unit of its own. A gain or a loss.
  \item \textbf{dBi} --- antenna gain relative to a theoretical isotropic
        radiator.
\end{tight}

\begin{worked}[The two rules that replace a calculator]
\textbf{The 3\,dB rule:} $+3$\,dB doubles power, $-3$\,dB halves it.\\
\textbf{The 10\,dB rule:} $+10$\,dB multiplies by ten, $-10$\,dB divides by ten.

So 20\,dBm is $10 \times 10 \times 1\,\text{mW} = 100\,$mW. And 23\,dBm is
double that again --- 200\,mW.

Going the other way: a client receiving $-70$\,dBm is receiving one ten-millionth
of a milliwatt. Radio receivers are extraordinarily sensitive, which is why small
amounts of noise matter so much.
\end{worked}

\subsection{Signal, noise and the number that actually matters}

\begin{longtable}{@{}L{2.8cm}L{5.5cm}L{6.7cm}@{}}
\toprule
\rowcolor{primary!15}
\textbf{Measure} & \textbf{What it is} & \textbf{Rule of thumb} \\
\midrule
\endhead
RSSI &
Received signal strength, in dBm. Always negative in practice &
$-67$\,dBm or better for voice and video; $-70$ for data;
below $-80$ is unusable \\
\rowcolor{lightbg}
Noise floor &
The background RF energy on the channel, in dBm &
Around $-92$\,dBm is clean. $-80$ is a problem \\
SNR &
The \emph{gap} between signal and noise, in dB &
25\,dB or better for voice; 20\,dB minimum for reliable data \\
\bottomrule
\end{longtable}

\begin{alertbox}[Signal strength alone tells you almost nothing]
A client at $-65$\,dBm sounds healthy. If the noise floor on that channel is
$-70$\,dBm, the SNR is 5\,dB and the link is unusable.

A client at $-75$\,dBm sounds marginal. With a noise floor of $-95$\,dBm the SNR
is 20\,dB and it works perfectly well.

\textbf{SNR is the number that decides whether a link works.} Users and helpdesk
tools report bars and RSSI; you should ask for the noise floor before drawing any
conclusion. This is the single most useful idea in the chapter and it comes up
again in Chapter~28.
\end{alertbox}

\subsection{What happens to a wave on the way}

\begin{longtable}{@{}L{3.2cm}L{5.6cm}L{6.2cm}@{}}
\toprule
\rowcolor{primary!15}
\textbf{Effect} & \textbf{What happens} & \textbf{Where you meet it} \\
\midrule
\endhead
Absorption &
Energy converted to heat by the material &
Walls, concrete, and \textbf{water} --- which means people. A room that works
when empty and fails when full \\
\rowcolor{lightbg}
Reflection &
The wave bounces off a smooth surface &
Metal, glass, lift shafts, warehouse racking. Causes multipath \\
Refraction &
The wave bends passing between materials &
Glass, water, atmospheric layers on long outdoor links \\
\rowcolor{lightbg}
Diffraction &
The wave bends around an obstacle, leaving a shadow behind it &
Pillars, corners. Why coverage dies just around a corner \\
Scattering &
The wave breaks up against an uneven surface or particles &
Rough plaster, chain-link fence, dust, heavy rain \\
\rowcolor{lightbg}
Free-space loss &
Signal weakens with distance regardless of obstacles &
Everywhere. Doubling the distance costs about 6\,dB \\
\bottomrule
\end{longtable}

\begin{conceptbox}[Antennas do not amplify]
An antenna with 9\,dBi of gain does not create energy. It \emph{shapes} the same
energy into a narrower pattern --- more in the direction you want, less
everywhere else.

\begin{tight}
  \item \textbf{Omnidirectional} --- radiates in a doughnut around the antenna.
        Ceiling APs in offices. Low gain, wide coverage.
  \item \textbf{Directional} (patch, Yagi, parabolic) --- radiates into a cone.
        Point-to-point links, long corridors, lecture theatres, warehouse aisles.
        High gain, narrow coverage.
\end{tight}

A common failure at scale: replacing omnidirectional antennas with high-gain ones
to ``fix'' coverage, and discovering the new pattern misses the floor directly
below the AP entirely.
\end{conceptbox}

\section{Security protocols}

\begin{longtable}{@{}L{1.9cm}L{2.9cm}L{3.6cm}L{6.1cm}@{}}
\toprule
\rowcolor{primary!15}
\textbf{} & \textbf{Encryption} & \textbf{Key exchange} & \textbf{Status} \\
\midrule
\endhead
WEP & RC4, static key & None worth the name &
\textbf{Broken since 2001.} Recoverable in minutes. Never deploy \\
\rowcolor{lightbg}
WPA & TKIP (RC4) & 4-way handshake &
Deprecated. A stopgap for WEP-era hardware \\
WPA2 & \textbf{AES-CCMP} & 4-way handshake, PSK or 802.1X &
The long-standing default. Vulnerable to offline dictionary attack on weak
pre-shared keys \\
\rowcolor{lightbg}
WPA3 & AES-GCMP &
\textbf{SAE} (Simultaneous Authentication of Equals) &
Current. Resists offline attack, adds forward secrecy, and adds OWE for open
networks \\
\bottomrule
\end{longtable}

\begin{examnote}[The two distinctions to get right]
\textbf{Personal against Enterprise} is about \emph{who} authenticates.

\begin{tight}
  \item \textbf{Personal (PSK)} --- one shared passphrase for everyone. Simple.
        Nobody is individually identified, and changing it means changing it on
        every device.
  \item \textbf{Enterprise (802.1X/EAP)} --- each user authenticates
        individually against a RADIUS server (\chref{aaa}). Revoke one account
        without touching anyone else. This is what a campus should run.
\end{tight}

\textbf{WPA2 against WPA3} is about \emph{how the key is agreed}. WPA2-Personal's
4-way handshake can be captured and attacked offline, so a weak passphrase falls.
WPA3 replaces it with SAE, where an attacker must guess against the live network
each time --- so offline attack is not possible, and each session has its own key
even if the passphrase is later disclosed.
\end{examnote}

\begin{conceptbox}[Open networks and OWE]
An open network with no passphrase sends everything in clear. \term{Opportunistic
Wireless Encryption}, part of the WPA3 family and marketed as ``Enhanced Open'',
encrypts the traffic anyway using an unauthenticated key agreement.

It stops passive eavesdropping on a guest network. It does \textbf{not}
authenticate the AP, so it does not stop somebody standing up a fake one. Worth
deploying on guest SSIDs; not worth mistaking for security.
\end{conceptbox}

\section{Causes of interference}

\subsection{Wi-Fi interfering with Wi-Fi}

\begin{longtable}{@{}L{4.2cm}L{5.2cm}L{5.6cm}@{}}
\toprule
\rowcolor{primary!15}
\textbf{} & \textbf{Co-channel (CCI)} & \textbf{Adjacent-channel (ACI)} \\
\midrule
\endhead
Cause & Two APs on the \emph{same} channel &
Two APs on \emph{partly overlapping} channels \\
\rowcolor{lightbg}
What happens & They hear each other and take turns politely &
They cannot decode each other, so they transmit over each other \\
Result & Reduced throughput. Everyone shares the airtime &
\textbf{Corruption and retransmission.} Much worse \\
\rowcolor{lightbg}
Verdict & Tolerable, and unavoidable at density &
To be eliminated. This is why the plan is 1, 6, 11 \\
\bottomrule
\end{longtable}

\begin{alertbox}[Counter-intuitive, and a favourite question]
\textbf{Co-channel interference is better than adjacent-channel interference.}

Two APs on channel 6 can hear one another, so CSMA/CA works: they defer to each
other and share the channel. Throughput halves; nothing is corrupted.

Two APs on channels 6 and 8 cannot decode one another --- each sees the other
only as raised noise. Neither defers, both transmit, and frames are destroyed and
retransmitted. Throughput collapses.

So a bad channel plan is worse than no channel plan. If you take one thing from
this section: \textbf{overlap, do not straddle.}
\end{alertbox}

\subsection{Everything else}

\begin{tight}
  \item \textbf{Microwave ovens} --- 2.4\,GHz, high power, intermittent. Classic
        symptom: wireless fails in the staff room at 12:30 every day.
  \item \textbf{Bluetooth and wireless headsets} --- 2.4\,GHz frequency hopping,
        low power, constant.
  \item \textbf{Cordless phones, video senders, wireless cameras, baby
        monitors} --- often 2.4\,GHz, often continuous, often invisible to
        Wi-Fi tools because they are not Wi-Fi.
  \item \textbf{Radar} --- 5\,GHz, and triggers DFS events rather than plain
        interference.
  \item \textbf{Too much of your own network} --- APs at full transmit power in a
        dense deployment interfere with each other. Reducing power often improves
        a busy site.
  \item \textbf{Physical obstruction} --- not interference strictly, but it
        presents identically to a user. Lift shafts, fire doors, foil-backed
        insulation, and fish tanks.
\end{tight}

\begin{conceptbox}[Non-Wi-Fi interference needs a spectrum analyser]
A Wi-Fi analyser on a laptop shows you Wi-Fi. A microwave oven does not appear in
it, because it transmits no frames --- it just raises the noise floor.

So a channel that looks empty in a Wi-Fi scan and performs terribly is the
signature of a non-Wi-Fi source. Finding it needs a spectrum analyser, which
looks at raw RF energy rather than at decoded frames. Knowing \emph{that the
distinction exists} is the examinable part.
\end{conceptbox}

\begin{pitfall}
\begin{tight}
  \item \textbf{Using any 2.4\,GHz channel other than 1, 6 or 11.} Creates ACI
        for two networks instead of sharing with one.
  \item \textbf{Bonding channels at 2.4\,GHz.} Consumes the band.
  \item \textbf{Reading RSSI without the noise floor.} SNR is the number that
        matters.
  \item \textbf{Turning transmit power up to fix coverage.} In a dense
        deployment this usually makes things worse, and it does nothing for the
        client's weaker transmitter on the return path.
  \item \textbf{Confusing SSID with BSSID.} A client roams between BSSIDs while
        the SSID never changes.
  \item \textbf{Assuming the AP controls roaming.} The client decides.
  \item \textbf{Treating WPA2-Personal as adequate for a campus.} Use Enterprise;
        a shared passphrase cannot be revoked for one person.
  \item \textbf{Expecting a Wi-Fi scanner to reveal a microwave oven.}
\end{tight}
\end{pitfall}

\begin{breakfix}[One lecture theatre is unusable when full and perfect when empty.]
A 200-seat theatre has four access points and tests clean every morning. During
lectures, clients associate but throughput collapses and voice calls fail. The
wired network behind the APs shows no errors and no congestion. A survey taken at
08:00 shows every seat at $-58$\,dBm or better.

% A listing nested inside a breakable box cannot itself be split, so
% without this tcolorbox forces it onto the page and reports an overfull
% vbox. Break the outer box early instead.
\needspace{8\baselineskip}
\begin{hostcmd}{Survey, taken during a full lecture}{wifi scan --- summary}
BSSID              SSID        Chan  Width  RSSI    Noise   SNR
00:11:22:aa:bb:01  CAMPUS      6     20     -61     -73      12
00:11:22:aa:bb:02  CAMPUS      6     20     -64     -73       9
00:11:22:aa:bb:03  CAMPUS      8     20     -59     -73      14
00:11:22:aa:bb:04  CAMPUS      11    20     -66     -73       7
\end{hostcmd}

\textbf{Diagnosis.} Three faults are visible in that one table, and they compound.

\textbf{First, the noise floor is $-73$\,dBm.} A clean environment is around
$-92$. Something is raising it by nearly 20\,dB, and the strong signal readings
are therefore worthless --- every SNR here is below the 20\,dB minimum for
reliable data.

\textbf{Second, the channel plan is wrong.} Four APs in one room are using
channels 6, 6, 8 and 11. The pair on channel 6 is co-channel interference, which
is merely wasteful. But \textbf{channel 8 overlaps both 6 and 11} --- it is
adjacent-channel interference against three other radios simultaneously, and that
is corruption rather than sharing.

\textbf{Third, the room is full of people.} Bodies are largely water, which
absorbs 2.4\,GHz strongly. This is why it works when empty: the morning survey
measured a room with no attenuation in it and no clients contending.

\textbf{Fix, in order of value.}

\begin{tight}
  \item Move the channel-8 radio to 1. That converts the worst interference in
        the room from ACI to CCI and costs nothing.
  \item Move as many clients as possible to 5\,GHz, where there is enough
        spectrum to give all four APs non-overlapping channels. In a dense room
        this is the real fix.
  \item Reduce transmit power so the four cells are smaller and clients associate
        with the nearest AP instead of the loudest.
  \item Find the noise source. Nothing in a Wi-Fi scan explains a $-73$\,dBm
        floor, so bring a spectrum analyser during a lecture.
\end{tight}

\textbf{And change the survey practice.} A survey of an empty room measures a room
that will never be used. Survey when occupied, or accept that the numbers are
optimistic by 10\,dB or more.
\end{breakfix}

\begin{lab}{A survey of somewhere you actually use}{Any laptop or phone with a Wi-Fi analyser}
\textbf{Platform note.} This topic has no configuration on the blueprint, so this
lab is measurement rather than building. Use any Wi-Fi analyser --- built-in
tools work: \cmd{netsh wlan show interfaces} on Windows, the option-click Wi-Fi
menu on macOS, \cmd{iw dev wlan0 scan} on Linux.

\textbf{Location.} A room you can visit twice, ideally once empty and once busy.

\begin{tasks}
  \item Record every visible SSID, its BSSIDs, band, channel and channel width.
        Draw the channel occupancy of 2.4\,GHz as a diagram like the one in this
        chapter.
  \item Identify every case of adjacent-channel interference in what you found.
        For each, name the channel that should have been used instead.
  \item Record RSSI \emph{and} noise for your own connection at five points:
        beside the AP, across the room, behind a wall, around a corner, and at
        the far edge of coverage. Calculate SNR at each.
  \item Plot RSSI against distance. Confirm the roughly 6\,dB loss per doubling
        of distance, and identify where your measurements depart from it and why.
  \item Find one place where signal is strong and SNR is poor. Explain what that
        tells you.
  \item Identify the security protocol in use on each SSID you can see. Note any
        still offering WPA or WEP.
  \item Move between two APs on the same SSID while watching the BSSID. Record
        the RSSI at which your client chose to roam --- and note that you did not
        choose it and the network did not either.
  \item \textbf{If you can:} repeat the RSSI and noise measurements with the room
        occupied. Compare, and quantify the absorption.
  \item \textbf{Extension.} Write the one-page recommendation you would give the
        building's network team, with the evidence for each point.
\end{tasks}
\end{lab}

\begin{checkpoint}
\begin{tightnum}
  \item Why does the 2.4\,GHz band have only three non-overlapping channels, and
        which are they?
  \item Which is worse, co-channel or adjacent-channel interference, and why?
  \item A client reports $-62$\,dBm and cannot hold a call. What is the next
        number you ask for, and what would make this normal?
  \item State the difference between WPA2-Personal and WPA2-Enterprise, and
        between WPA2 and WPA3.
  \item Give two reasons a room might work when empty and fail when full.
  \item What kind of interference is invisible to a Wi-Fi analyser, and what tool
        finds it?
\end{tightnum}
\end{checkpoint}

\begin{chaptersummary}
\begin{tight}
  \item Wireless is half duplex, shared, and uses collision \emph{avoidance}. The
        advertised rate is shared between everyone on the channel.
  \item SSID is the name, BSSID is the radio's MAC. The \emph{client} decides
        when to roam.
  \item 2.4\,GHz: three non-overlapping channels, 1, 6 and 11. 5\,GHz: around 25,
        some DFS. 6\,GHz: 59, newest clients only.
  \item Channel bonding doubles throughput, halves the channel count and raises
        the noise floor. Never at 2.4\,GHz.
  \item $+3$\,dB doubles power, $+10$\,dB multiplies by ten. 0\,dBm is 1\,mW.
  \item \textbf{SNR decides whether a link works}, not RSSI. Aim for 25\,dB for
        voice, 20\,dB minimum for data.
  \item Radio waves are absorbed, reflected, refracted, diffracted and scattered;
        free-space loss costs about 6\,dB per doubling of distance.
  \item WEP is broken, WPA deprecated, WPA2 uses AES-CCMP with a 4-way handshake,
        WPA3 uses SAE and resists offline attack. Personal shares a passphrase;
        Enterprise authenticates each user via RADIUS.
  \item Co-channel interference shares politely; adjacent-channel interference
        corrupts. A bad plan is worse than no plan.
\end{tight}
\end{chaptersummary}

\begin{reviewq}
  \item \textbf{(MCQ)} Which set of 2.4\,GHz channels does not overlap?
        \begin{tight}
          \item[(a)] 1, 5, 9 \quad (b)~1, 6, 11 \quad (c)~2, 7, 12
          \quad (d)~1, 4, 8
        \end{tight}
  \item \textbf{(MCQ)} A client shows RSSI $-68$\,dBm and a noise floor of
        $-88$\,dBm. What is the SNR, and is it adequate for data?
        \begin{tight}
          \item[(a)] 20\,dB, adequate \quad (b)~20\,dB, inadequate
          \item[(c)] 156\,dB, adequate \quad (d)~$-20$\,dB, inadequate
        \end{tight}
  \item \textbf{(MCQ)} Which WPA3 mechanism replaces the WPA2 4-way handshake and
        prevents offline dictionary attack?
        \begin{tight}
          \item[(a)] TKIP \quad (b)~CCMP \quad (c)~SAE \quad (d)~OWE
        \end{tight}
  \item \textbf{(Short)} Explain why two APs on channel 6 perform better than two
        APs on channels 6 and 8.
  \item \textbf{(Short)} An outdoor link works in dry weather and fails in heavy
        rain. Name the two RF effects responsible and explain each.
  \item \textbf{(Applied)} You are given a floor with six APs, all on 2.4\,GHz, on
        channels 1, 3, 6, 8, 11 and 11. Propose a corrected plan, justify each
        change, and state the one further recommendation that would help more
        than any channel change.
  \item \textbf{(Applied)} A site reports that wireless drops for a few seconds at
        unpredictable times on 5\,GHz only, and never on 2.4\,GHz. Give the most
        likely cause, explain the mechanism, and state how you would confirm it.
\end{reviewq}
